\documentclass[11pt]{article}

\usepackage[margin=1in]{geometry}
\usepackage{amsmath}
\usepackage{amssymb}
\usepackage{booktabs}
\usepackage{array}
\usepackage{longtable}
\usepackage{enumitem}
\usepackage{xcolor}
\usepackage{siunitx}
\usepackage{hyperref}
\hypersetup{
  colorlinks=true,
  linkcolor=black,
  urlcolor=blue,
  citecolor=black
}

\newcommand{\ppm}{\,\text{ppm}}
\newcommand{\Da}{\,\text{Da}}
\newcommand{\Th}{\,\text{Th}}
\newcommand{\Cthirteen}{\textsuperscript{13}C}
\newcommand{\mz}{\ensuremath{m/z}}
\newcommand{\zeq}[1]{\ensuremath{z{=}#1}}

\title{Failure Modes in Untargeted MS1 Feature-Detection Algorithms}
\author{Alex Solivais}
\date{\today}

\begin{document}
\maketitle

\begin{abstract}
This chapter is an empirical study of \emph{why} an untargeted (de-novo, no-identification)
MS1 feature-detection pipeline fails to recover peptide features, and of what fixes
each failure demands. The vehicle is a pure-Rust reimplementation of FlashLFQ extended
with a three-stage untargeted pipeline (detect $\rightarrow$ refine $\rightarrow$ resolve),
validated by matching its output against FlashLFQ's identification-driven quantified peaks
on a carbonic-anhydrase tryptic HCD run acquired on a Thermo Orbitrap Fusion Lumos. Over one
extended investigation, reference recall rose from $589/620$ ($95.0\%$) to $602/620$ ($97.1\%$)
and charge-matched recall from $564$ ($91.9\%$) to $584$ ($94.2\%$) by diagnosing and repairing
a sequence of distinct failure modes. Rather than present only the working pipeline, this chapter
organises those failures into a taxonomy by pipeline stage and by root mechanism, and for each
gives symptom, mechanism, concrete evidence (named peptide with masses, \mz, retention time, and
fit scores), the fix, and the generalisable lesson. It devotes a dedicated section to
\emph{negative} results --- fixes that regressed recall and were therefore reverted --- because
those are as instructive as the successes and prevent repeated effort. It closes with design
principles for isotope-envelope feature detectors and a candid accounting of the misses that
remain. Where a claim rests on a single traced example, this is stated explicitly.
\end{abstract}

\tableofcontents

\section{Scope, pipeline, and validation methodology}
\label{sec:scope}

\subsection{The pipeline under study}
The system is a de-novo MS1 feature detector: it receives centroided MS1 scans and must
produce a list of peptide-level features (monoisotopic neutral mass, charge state(s),
retention-time (RT) extent, intensity) with \emph{no} peptide identifications as input. It
has three stages.

\begin{enumerate}[leftmargin=1.6em]
  \item \textbf{Detect} (\texttt{trace\_kernel.rs}). A sparse matched-filter / trace-kernel
  detector. Every centroid is a candidate seed; seeds are processed tallest-first. For each
  seed, a charge is decided by scoring an isotope-comb $\times$ RT-Gaussian hypothesis over
  \zeq{z} for $z \in [z_{\min}, z_{\max}]$ and keeping the highest response (\emph{cross-charge
  non-maximum suppression}, cross-$z$ NMS). The winning feature's peaks are then \emph{claimed}
  --- removed from further seeding --- so overlapping harmonics cannot re-fire. The default
  scorer is \texttt{RawSum}, an un-normalised $\sum_k \sum_s w_k\, g_s\, I_{k,s}$ over comb
  teeth $k$ and RT scans $s$. A coverage-target stop and a chromatographic-persistence gate
  ($\geq 2$ distinct scans) bound the feature population.
  \item \textbf{Refine} (\texttt{feature\_refinement.rs}, \texttt{isotope\_shift\_decon.rs}).
  Each detected feature is deconvolved on an averaged apex composite ($\leq 3$ scans, apex $\pm 1$)
  and/or the single apex scan. A FlashLFQ-style \Cthirteen{} isotope-shift search ($-1/0/+1$, widened
  to $-2\ldots+2$ at high charge) places the monoisotope; an \emph{envelope-fit cosine} metric
  selects charge and monoisotope. A high-charge \emph{mono walk-back} corrects heavy-peptide
  off-by-one-too-high placements.
  \item \textbf{Resolve} (\texttt{feature\_refinement::resolve\_charge\_state\_consensus}).
  Refined features are knitted into peptide-level features by single-linkage grouping on
  co-elution plus off-by-one-aware mass agreement, then a cross-charge mass vote settles the
  consensus monoisotopic mass.
\end{enumerate}

\subsection{Validation data and metric}
\label{sec:validation-data}
The benchmark is a carbonic-anhydrase (``CA'') tryptic digest, HCD fragmentation, $10$-minute
gradient, acquired on a Thermo Lumos (\num{2689} MS1 scans, uncalibrated \texttt{.raw}). The
reference is FlashLFQ's \texttt{AllQuantifiedPeaks.tsv} from a wide-tolerance
MetaMorpheus/GPTMD search ($634$ rows; $\sim620$ quantifiable peaks after filtering MBR/decoy
rows with non-numeric RT/intensity). \textbf{Recall} is the fraction of reference peaks
rediscovered by the untargeted output within $\pm 20\ppm$ mass and $\pm 0.3$ min RT;
\textbf{charge-matched recall} additionally requires the charge state to agree.

Two methodological cautions are load-bearing throughout this chapter.

\begin{itemize}[leftmargin=1.4em]
  \item \textbf{A systematic \boldmath$\sim+10\ppm$ calibration offset.} The reference
  \texttt{Peptide Monoisotopic Mass} is the \emph{theoretical} peptide mass; the raw file is
  \emph{uncalibrated}. Untargeted masses therefore run systematically $\sim +10\ppm$ high
  against the reference. This is instrument calibration, not a detector error. Concretely,
  EFLDVIK's theoretical reference mass is $862.480\Da$ while the untargeted pipeline resolves
  it at $862.488\Da$ --- a $\sim 9\ppm$ offset, not a mono-placement error. Matches are
  accepted within an integer-\Cthirteen{} offset plus residual ppm.
  \item \textbf{Compare on observed \mz, not theoretical mass.} The reference \texttt{Peak MZ}
  is the \emph{most-abundant} isotope \mz, not the monoisotopic \mz; for $z \geq 2$ it sits
  $\sim 1$ \Cthirteen{} above the mono ($\approx +0.5\Th$ at \zeq{2}). Observed-vs-observed \mz
  agreement is dead-on ($\pm 1$--$2\ppm$); the residual $\sim +10\ppm$ appears only in
  theoretical-vs-observed comparisons.
\end{itemize}

\subsection{How to read this chapter}
The failure modes are organised by pipeline stage (Sections~\ref{sec:mono}--\ref{sec:detector}),
because the stage locates the fix, but the deeper organising axis is \emph{root mechanism}: most
failures reduce to one of four mechanisms --- (i) a scoring window that cannot see signal on one
side of the hypothesis; (ii) harmonic ($z \leftrightarrow 2z$) ambiguity inherent to isotope-comb
matching; (iii) anchoring a deconvolution on the wrong (stronger, foreign) peak; and (iv)
aggregating cross-charge evidence by quantity or intensity rather than by fit quality.
Table~\ref{tab:summary} maps every documented failure to its stage, mechanism, traced peptide,
and fix; Table~\ref{tab:progression} records the recall progression across the fixes.

\section{Monoisotope placement}
\label{sec:mono}

\subsection{Heavy-peptide off-by-one-too-high, and the window blind spot}
\label{sec:walkback}

\paragraph{Symptom.} Heavy, highly-charged peptides are recovered at the correct charge and
approximate mass but with the monoisotope placed one or two \Cthirteen{} units too high, so the
neutral mass is $\sim 1$--$2\Da$ heavy and the reference match fails.

\paragraph{Mechanism.} For a heavy peptide the isotopic envelope \emph{mode} (its tallest tooth)
sits one or more \Cthirteen{} above the true monoisotope. The detector seeds on the tallest peak;
the shift search, when its anchor or slice differs from the ideal, can then place the mono on a
peak one or two \Cthirteen{} above the true one. Crucially, the envelope-fit cosine metric that
scores placements uses a window that begins at $\text{mono} - 0.5\cdot\Delta$ (where $\Delta$ is
the isotope spacing). A too-high placement's window therefore \emph{never sees} the strong true-mono
peak sitting just below it: the unexplained signal beneath the hypothesis is off-grid, so the
wrong placement scores as cleanly as the right one. This is the \emph{window blind spot}: a fit or
correlation window that does not extend below the monoisotope is structurally incapable of
detecting an off-by-one-too-high.

\paragraph{Evidence.} Traced on \textbf{ECCHGDLLECADDRADLAK} (reference monoisotopic mass
$2246.9355\Da$, \zeq{4}, RT $12.153$~min; mono \mz\ $562.741$, $+1$ at $562.997$, $+2$ at
$563.243$). The detector seeded the mono on the $+2$ isotope ($562.247$, mass $\approx 2248.96$).
A placement one \Cthirteen{} too high (mass $2247.9578$) scores cosine $0.934$ --- an apparently
acceptable fit --- precisely because the strong $562.746$ peak ($\approx\!8.0\times10^{6}$
intensity) lies just below the window and is invisible to the metric.

\paragraph{Fix and why it works.} \texttt{walkback\_mono\_high\_charge} re-scores the chosen mono
against its $1$- and $2$-\Cthirteen-lower alternatives and adopts the \emph{lowest-mass} hypothesis
whose fit is within a small margin ($\texttt{WALKBACK\_FIT\_MARGIN} = 0.01$) of the best. The lower
alternatives' windows \emph{do} include the peaks beneath the original placement, so they can finally
see the unexplained signal that condemned it. The check is self-protecting: a correctly-placed mono's
lower alternatives predict a leading tooth where nothing is observed, so they fit far worse and a good
mono is never walked back. It is gated on mass $\geq 2200\Da$ and $|z| \geq 3$ --- the heavy-peptide
regime where the mode sits well above the mono. On ECC the walk-back recovers $2246.953\Da$ (a dead-on
reference match within the calibration offset).

\paragraph{Generalisable insight.} A model-selection metric must be able to see signal on
\emph{both} sides of the hypothesis it scores. A one-sided window silently converts an entire
directional error class (here, mono-too-high) into a blind spot; the fix is either to widen the
window (but see Section~\ref{sec:neg-window} for why na\"ive widening regressed) or to add an
explicit, gated re-scoring against lower-mass alternatives whose windows \emph{do} reach the
missing evidence.

\section{Charge assignment: harmonic ambiguity}
\label{sec:charge}

Isotope-comb matching has an inherent degeneracy: the teeth of a charge-\zeq{z} envelope
(spacing $\Delta = 1.00335\Da / z$) are a \emph{subset} of the charge-\zeq{2z} grid (spacing
$\Delta/2$). A \zeq{z} hypothesis and a \zeq{2z} hypothesis can therefore both ``explain'' the
same peaks. Which way the ambiguity bites depends on what fills --- or fails to fill --- the
intervening teeth.

\subsection{Charge-halving (real \boldmath\zeq{2} labelled \zeq{1})}
\label{sec:halving}

\paragraph{Symptom.} A genuine light \zeq{2} peptide is detected but assigned \zeq{1}, halving its
neutral mass, so it never matches the true-mass reference. This class accounted for $\sim 10$
peptides (light \zeq{2}, $788$--$907\Da$, \mz\ $\sim 395$--$455$).

\paragraph{Mechanism.} A \zeq{2} envelope's \emph{even} isotopes (\mz\ $432.25, 433.25, 434.25,
\ldots$) themselves form a valid $1.0$-\Th\ \zeq{1} comb. The \texttt{RawSum} charge scorer
(Eq.~\ref{eq:rawsum-ch}) rewards matched intensity but never penalises the strong \emph{intervening}
\zeq{2} peaks (the $+1$ at $432.75$) that a \zeq{1} hypothesis leaves unexplained. \zeq{2} explains
every observed peak; \zeq{1} explains every \emph{other} one and ignores the rest --- yet \zeq{1}
out-scores \zeq{2}. Refine and resolve then faithfully carry the halved mass.
\begin{equation}
\texttt{RawSum}(z) \;=\; \sum_{k}\sum_{s} w_k\, g_s\, I_{k,s},
\label{eq:rawsum-ch}
\end{equation}
un-normalised and never decremented, is blind to both missing predicted teeth and unexplained
observed peaks.

\paragraph{Evidence.} Traced on \textbf{EFLDVIK} (reference mass $862.480\Da$ theoretical, true
\zeq{2}, RT $14.70$~min; most-abundant \mz\ $432.251$). The strongest feature at that \mz\
($432.28$, intensity $8.7\times10^{8}$) was seeded and assigned \zeq{1} $\rightarrow$ mono
$431.28\Da$ = half the true $862.48\Da$. Raising the coverage target ($0.90\rightarrow 0.95$;
\num{96000}$\rightarrow$\num{143000} features) recovered \emph{zero} additional reference peaks and
left the miss set byte-identical --- proof that this is a charge-scoring failure, not a detection-depth
one. The envelope-fit cosine cleanly ranks the charges: \zeq{2} $0.984 >$ \zeq{1} $0.905 >$ \zeq{3}
$0.718$ (see the companion note \emph{An Envelope-Fit Metric for Charge-State and Monoisotope
Selection} for the full derivation and the completeness/explained-fraction decomposition).

\paragraph{Fix and why it works.} Charge \emph{re-selection} in refine
(\texttt{best\_charge\_by\_fit}) re-scores the candidate charges $\{z/2, z, 2z\}$ by the
envelope-fit cosine, which \emph{does} penalise the intervening peaks the \zeq{1} comb leaves
unexplained (they enter the cosine as $(t=0, o=I)$ slots that inflate $\lVert o \rVert$ with no
numerator contribution). The metric places the mono where the whole envelope is explained, not
where a subset comb happens to land. This made detector-anchored shift decon $+$ envelope-fit
recharge the pipeline default and took recall $89.5\% \rightarrow 95.0\%$ (charge $85.3\%
\rightarrow 91.0\%$); the ``absent-miss'' class fell $13 \rightarrow 3$. EFLDVIK now resolves at
$862.488\Da$ / \zeq{2}.

\paragraph{Generalisable insight.} A charge scorer used for cross-$z$ NMS must reward
\emph{explaining the intervening teeth}, not just matched intensity; equivalently, it must
penalise the signal a lower-charge subset comb ignores. \texttt{RawSum}'s un-normalised matched
filter cannot do this by construction, which is the structural motivation for the envelope-fit
cosine.

\subsection{Charge-doubling (real \boldmath\zeq{z} labelled \zeq{2z})}
\label{sec:doubling}

\paragraph{Symptom.} The mirror image: in crowded windows, charge re-selection promotes a genuine
\zeq{z} feature to \zeq{2z}, doubling its mass and losing the peak. This is the danger the
recharge machinery itself introduces, and it was widespread, not confined to the traced case.

\paragraph{Mechanism.} Because the \zeq{z} teeth are a subset of the \zeq{2z} grid, and in a
crowded window co-eluting \emph{interferents} happen to fill the intervening \zeq{2z} teeth, the
\zeq{2z} hypothesis ``explains more'' and wins a strict argmax over $\{z/2, z, 2z\}$. A relative
margin alone cannot save the true charge here: the correct \zeq{z} envelope may genuinely fit
\emph{poorly} (weak, overlapped), so there is no relative gap to defend it with.

\paragraph{Evidence.} Traced on the sodium adduct \textbf{HVGDLGNVTAD[Na]K} (reference mass
$1246.5918\Da$, true \zeq{2}, RT $11.238$~min, weak at $2.75\times10^{6}$; mono \mz\ $624.309$).
The detector had \zeq{2} right, but strict-argmax recharge doubled it to \zeq{4} (mass $2492.20$).
The measured fits are \zeq{1} $0.235$, \zeq{2} $0.258$, \zeq{4} $0.689$: the correct \zeq{2}
envelope \emph{genuinely fits poorly}, and the spurious \zeq{4} ``win'' is only a mediocre cosine
($\sim0.69$) earned by absorbing interferents on the denser grid. No margin over the true charge
exists to protect it.

\paragraph{Fix and why it works.} \texttt{best\_charge\_by\_fit} keeps the detector's own charge
(\texttt{prefer\_charge}) unless an alternative \emph{both} beats it by
$\texttt{RECHARGE\_PREFER\_MARGIN}=0.05$ \emph{and} itself clears an \emph{absolute} floor
$\texttt{RECHARGE\_MIN\_FIT}=0.85$. The absolute floor is the key discriminator: a spurious
harmonic win is only $\sim0.69$, well under the floor, whereas a genuine re-charge (e.g. the light
\zeq{2} charge-halving recovery of Section~\ref{sec:halving}) fits cleanly at $\sim0.98$ and still
fires. HVG now resolves at $1246.60\Da$ / \zeq{2}. Blocking spurious doubling took recall
$96.6\% \rightarrow 97.1\%$ and charge-matched $92\% \rightarrow 94\%$ --- a dataset-wide effect.

\paragraph{Generalisable insight.} Harmonic ($z \leftrightarrow 2z$) ambiguity is intrinsic to
isotope-comb matching and cannot be resolved by relative fit comparison alone. It requires (a)
\emph{loyalty} to a prior (here the detector's own charge, itself derived from independent
evidence) and (b) an \emph{absolute} quality gate, so that a charge is only overridden by an
alternative that is good in absolute terms, not merely better than a poor incumbent. Relative
margins are necessary but not sufficient when the true hypothesis can itself fit badly.

\section{Anchor selection}
\label{sec:anchor}

\subsection{The distant grab}
\label{sec:distant-grab}

\paragraph{Symptom.} A feature's deconvolution places the monoisotope several isotopes away, on a
different species, drifting the mass by multiple \Cthirteen{} units.

\paragraph{Mechanism.} When deconvolution anchors on the \emph{tallest peak in a wide window}, and
that window contains a stronger co-eluting \emph{different} species several isotopes above the true
feature, the anchor lands on the neighbour and the whole averagine template is placed relative to
it --- putting the mono on the interloper. This ``distant grab'' was the root of $\sim78\%$ of the
disagreements in a four-way (classic$\times$shift $\times$ composite$\times$apex) decon comparator.

\paragraph{Fix and why it works.} Anchor the shift views on the \emph{detector's own} most-abundant
claimed peak (\texttt{four\_way\_decon\_detector\_anchor}; the tallest of \texttt{feature.peaks}),
so the deconvolution cannot grab a foreign peak at all. A neighbour-grid gating variant
(\texttt{shift\_decon\_gated}, which excludes peaks on a co-eluting \emph{detected} neighbour's
isotope grid unless they are on this feature's own grid) was also implemented, but it only reduced
disagreement to $60.8\%$ because \emph{undetected} neighbours slip through; detector-anchoring
collapsed disagreement $64.2\% \rightarrow 45.4\%$ (unanimous $35.8\% \rightarrow 54.6\%$) and won.

\paragraph{Averagine keyed on the most-intense mass, not the monoisotope.} A related design choice
underlies the anchor: the averagine template is keyed by the observed \emph{most-intense} (tallest)
peak's mass, not by a monoisotope guess. The tallest peak is the highest-SNR, most accurately
centroided, hardest-to-misassign point in the envelope; the monoisotope --- especially for heavy
peptides whose mode sits above it --- is exactly the fragile quantity being recovered. Keying the
template on the mono would be circular. The template's mode tooth is pinned at the anchor and the
mono falls out as $\text{anchor} - \text{mode}\cdot\Delta$, corrected by the shift search. This
mirrors the classic deconvolution's use of the observed most-intense peak as reference.

\paragraph{Generalisable insight.} Anchor a deconvolution on the most robust observable, not on
the quantity you are trying to estimate. When a window may contain foreign species, the safest
anchor is one an upstream stage already committed to (the detector's claimed peak), because it
cannot drift onto a neighbour.

\subsection{The single apex scan beats the averaged composite}
\label{sec:apex}
A ground-truth cut (per-feature verdicts joined to the FlashLFQ peaks) showed that when the
apex-composite and single-apex views disagree, the \emph{single apex scan} is correct $70\%$ of the
time (classic decon) versus $16\%$ for the composite; per-view accuracy is classic-apex $90.2\%$,
shift-apex $87.7\%$, shift-composite $86.7\%$, classic-composite $83.9\%$. Averaging even a small
window pulls in co-eluting interference; the apex is a higher-SNR snapshot. This motivates the
domain rule that the composite must stay small (Section~\ref{sec:open}) and the pipeline's default
to the shift-apex view. \emph{Caveat:} these accuracies are from one dataset with narrow ($\sim1.9$~s
FWHM) peaks; on longer gradients a slightly wider composite may be preferable.

\section{Cross-charge resolution and knitting}
\label{sec:resolve}

The resolve stage knits refined features into peptide-level features. Its two sub-decisions ---
\emph{which} features to merge, and \emph{which} mass to adopt for the merged group --- each have a
characteristic failure.

\subsection{Merging separately-eluting species as false off-by-ones}
\label{sec:knit}

\paragraph{Symptom.} A correctly-placed feature is dragged onto a wrong mass because it is knitted
with a \emph{different} co-eluting species $\sim 1\Da$ away, which then out-votes it.

\paragraph{Mechanism.} \texttt{features\_link} bridged any two features within $\pm 2$ \Cthirteen{}
in mass and within the apex-RT window, with \emph{no co-elution check}. A different peptide
$\sim 1\Da$ away that merely drifted inside the $0.1$-min window was knitted as a false off-by-one;
being stronger, it dragged the consensus mass onto itself.

\paragraph{Evidence (must split).} For \textbf{ECCHGDLLECADDRADLAK} (\zeq{4}, apex $12.153$~min,
narrow elution window $[12.153, 12.169]$), the corrected \zeq{4} feature was wrongly merged with a
\emph{different} co-eluting \zeq{2}/\zeq{3} species at $\sim 2247.94\Da$ (apex $12.252$~min, outside
\zeq{4}'s window). The stronger interloper out-voted the corrected \zeq{4} and the true mono was
lost. Note additionally that at \zeq{3}/\zeq{2} ECC's true mono peak ($749.986$ / $1124.475$~\mz)
is \emph{genuinely absent} from the centroids (its envelope runs $+1, +2, +3, +4$), so those
charges legitimately fit $+1$; only \zeq{4} sees the mono --- see Section~\ref{sec:datagaps}.

\paragraph{Evidence (must stay knitted --- counter-example).}
\textbf{DSCQGDSGGPVVCNGQLQGIVSWGYGCAQK} (reference mass $3183.3808\Da$, \zeq{3}, RT $15.414$~min)
genuinely co-elutes at \zeq{3}/\zeq{4}/\zeq{5} (all apex $15.414$~min). These are the \emph{same}
peptide at different charges and must remain knitted; the fix must not break them.

\paragraph{Fix and why it works.} Off-by-one links ($k \neq 0$) are gated on tight co-elution
(\texttt{features\_coelute\_tightly}): each feature's apex must fall inside the other's
$[\text{start\_rt}, \text{end\_rt}]$ window (padded by $\pm 0.6$~s). Pure same-mass cross-charge
links ($k = 0$, the same peptide at another charge) are untouched, so DSC's genuinely co-eluting
charges still knit, while ECC's \zeq{4} splits from the \zeq{2}/\zeq{3} interloper whose apex lies
outside \zeq{4}'s narrow window. This was the single largest lever of the session: $+9$ recall
(part of $589 \rightarrow 598$), not just the traced peptide.

\paragraph{Generalisable insight.} Chromatographic co-elution is a discriminator the mass domain
alone does not have. Two species $\sim 1\Da$ apart are indistinguishable to an off-by-one mass test,
but their elution profiles are not: requiring mutual apex-in-window co-elution before treating a
$\sim 1\Da$ mass difference as an off-by-one uses the orthogonal RT axis to separate real
mono-placement disagreements (same peptide) from coincidental mass neighbours (different peptides).

\subsection{Consensus voting by count or intensity rather than fit}
\label{sec:vote}

\paragraph{Symptom.} The knitted group settles on the wrong monoisotopic mass even though a member
placed it correctly.

\paragraph{Mechanism, two variants.}
\begin{enumerate}[label=(\alph*),leftmargin=1.8em]
  \item \textbf{Count.} Two off-by-one siblings out-vote one correct placement when the group votes
  by candidate count.
  \item \textbf{Intensity.} A contaminated-but-stronger charge out-votes a clean but weaker one when
  ties are broken by summed intensity.
\end{enumerate}

\paragraph{Evidence (intensity variant).} \textbf{AAVTAFWGK} (reference mass $992.508\Da$;
FlashLFQ quantified it at \zeq{1} although the precursor was \zeq{2}; RT $16.31$~min; \zeq{1} mono
\mz\ $993.515$). It is seen at \zeq{1} (envelope contaminated by co-eluting interferents; refine
mis-places the mono $+1$, envelope-fit $0.39$) and at \zeq{2} (clean; correct mono, fit $0.99$).
Both are lone singletons (support $1$), so an intensity-only tiebreak took the stronger-but-wrong
\zeq{1} $+1$ mass and lost the true mono.

\paragraph{Fix and why it works.} \texttt{resolve\_mass\_by\_cross\_charge} ranks mass clusters by
the tuple $(\text{support},\, \text{count},\, \text{best decon\_score},\, \text{summed intensity})$.
Distinct-charge \emph{support} stays the \emph{primary} key --- a genuine multi-charge agreement must
still outrank a lone high-fit placement --- but the envelope-fit is inserted as a tiebreak
\emph{before} intensity, so when support ties the clean placement wins. AAVTAFWGK now resolves at
$992.517\Da$ / \zeq{1};\zeq{2}. Recall $598 \rightarrow 599$. Critically, promoting summed fit to the
\emph{primary} key regressed recall (Section~\ref{sec:neg-fitprimary}); fit works only as a tiebreak
after support.

\paragraph{Generalisable insight --- cross-charge redundancy is the central resource.} When one
charge state is ambiguous, contaminated, or missing its monoisotope, another charge of the same
peptide often resolves it. But this redundancy pays off \emph{only if resolution weights quality
over quantity and intensity}. A stronger contaminated charge and two off-by-one siblings are both
``more'' evidence pointing the wrong way; fit quality is what distinguishes a clean placement from a
loud wrong one. This is arguably the chapter's central positive insight: exploit cross-charge
redundancy, but weight it by fit, with independent multi-charge agreement (support) as the primary
anchor so a single high-fit outlier cannot override a genuine consensus.

\section{Data-limited and inherent-ambiguity cases}
\label{sec:datagaps}

Not every off-by-one is an algorithm bug. Two classes are genuine data limitations, and honesty
about them bounds what any algorithm can achieve.

\paragraph{Missing monoisotope at some charges.} For some peptides the true monoisotope peak is
simply \emph{absent} from the centroids at certain charges --- contamination or peak-picking drops
it --- while present at others. \textbf{ECC} again: at \zeq{2}/\zeq{3} the mono \mz\ ($1124.475$ /
$749.986$) is not in the centroids at all (the observed envelope begins at $+1$), so the envelope
legitimately starts one \Cthirteen{} above the mono and the fit \emph{correctly} prefers the $+1$
placement at those charges. The ``error'' is a real data gap. It is resolvable only by trusting the
charge that \emph{does} see the mono (\zeq{4} here) --- which is exactly why cross-charge resolution
(Section~\ref{sec:vote}) with quality weighting, not per-charge fit, is the right backstop.

\paragraph{Contaminated envelopes in crowded regions.} Even the \emph{correct} charge can fit
poorly when the window is crowded: AAVTAFWGK \zeq{1} fits $0.39$; HVG \zeq{2} fits $0.258$. In such
regions absolute cosine values are low across all hypotheses and relative fit comparisons become
unreliable --- the reason Section~\ref{sec:doubling}'s absolute floor and Section~\ref{sec:vote}'s
support-primary ranking are needed rather than pure argmax-of-fit.

\paragraph{Insight.} Distinguish genuine algorithm errors from data-limited inherent ambiguity.
Some ``misses'' are cases where the monoisotope is not in the data at a given charge, or where the
reference itself may be questionable. Chasing these with ever-more-aggressive placement heuristics
risks regressing the cases that \emph{are} recoverable (see the negative results). The correct
response is architectural: rely on the charge that sees the clean signal, via quality-weighted
cross-charge consensus.

\section{Detector-level artifacts}
\label{sec:detector}

\paragraph{Greedy tallest-first claiming.} The detector claims peaks tallest-first, and each peak
is owned by exactly one feature (a claim partition). This is what makes cross-$z$ NMS work
(a claimed harmonic cannot re-fire), but it creates artifacts that interact with refinement: in
dense regions a real isotope tooth can be claimed first by a marginally-taller \emph{spurious}
neighbour, so the claim partition is noisy. This directly explains why refine-stage peak censoring
regressed recall (Section~\ref{sec:neg-censor}): deleting ``already-claimed'' peaks removes real
teeth that the noisy partition mis-assigned.

\paragraph{The \texttt{RawSum} charge score.} As formalised in Eq.~\ref{eq:rawsum-ch} and
Section~\ref{sec:halving}, the default detector charge score is an un-normalised
$\sum_k\sum_s w_k g_s I_{k,s}$ that neither penalises missing predicted teeth nor unexplained
observed peaks. It is the upstream root of the charge-halving class and the structural motivation
for the envelope-fit cosine downstream. A normalized matched-filter variant
(\texttt{ScoreModel::NormalizedNoiseFloor}) was implemented and measured; it gave $+1.4\%$ recall
but $-0.8\%$ charge and so failed its ``charge up or flat'' bar, and \texttt{RawSum} remains the
detector default --- the charge fix was moved downstream into refine's envelope-fit recharge, which
is where the cosine can see the whole envelope on a clean apex view.

\paragraph{Insight.} A greedy claim partition is an implementation convenience, not ground truth;
downstream stages must not treat ``claimed by another feature'' as ``not part of my envelope.''
And the natural place to fix a charge error is wherever the richest evidence lives --- here,
refine, where a high-SNR apex composite and a completeness-aware metric are both available.

\section{Negative results}
\label{sec:negative}

These experiments regressed recall and were reverted. They are recorded because each rules out a
plausible direction and explains \emph{why} the intuition failed.

\subsection{Widening the envelope-fit window below the mono}
\label{sec:neg-window}
To attack the window blind spot (Section~\ref{sec:walkback}) directly, the envelope-fit window's
lower reach was extended from $\text{mono}-0.5\cdot\Delta$ to $\text{mono}-2.5\cdot\Delta$ so the
metric could see the peak below a too-high mono. This regressed recall both globally (to $586$) and
when scoped only to the walk-back ($587$). \textbf{Why it fails:} a wider lower window admits
\emph{legitimate} below-mono signal --- the $+1$/$+2$ teeth of a genuine lower-mass co-eluting
species, or noise --- as spurious ``unexplained'' penalty against \emph{correct} placements,
depressing good cosines more than it rescues bad ones. The targeted, self-protecting walk-back (which
only \emph{prefers} a lower mono when it fits at least as well) succeeds precisely because it does
not globally change what counts as in-window.

\subsection{Fit-weighted consensus as the primary ranking key}
\label{sec:neg-fitprimary}
Ranking mass clusters by \emph{summed} \texttt{decon\_score} as the \emph{primary} key (instead of
distinct-charge support) regressed recall. \textbf{Why it fails:} two off-by-one siblings still
out-\emph{sum} one correct member, so summing fit reproduces the count failure it was meant to cure.
This is the crucial contrast with Section~\ref{sec:vote}'s working fix, where fit is only a
\emph{tiebreak after support}: independent multi-charge agreement must remain the primary anchor, or
a pair of concordant-but-wrong placements overrides a lone correct one.

\subsection{Over-broad walk-back gating}
\label{sec:neg-walkback}
The mono walk-back was tried at all charges, with only a $2000\Da$ mass floor, and with an
$\text{argmax(template)} \geq 1$ mode gate ($\sim1500\Da$). All regressed (to $587$--$588$).
\textbf{Why it fails:} at low charge / mid mass the fit metric's window cannot see below the mono
(the same blind spot), so it scores a $+1$ placement above the true one and the walk-back
\emph{over-corrects}, walking a correct mid-mass low-charge mono down to a wrong lower mass. The
gate must key on the physical driver of the error --- high mass \emph{and} high charge, where the
mode genuinely sits above the mono --- hence $\geq 2200\Da$ and $|z| \geq 3$.

\subsection{Refine-stage subtractive censoring}
\label{sec:neg-censor}
Removing peaks claimed by \emph{other} features from a feature's composite before deconvolution
(to clean chimeric windows) regressed recall substantially: baseline $88.2\%$ / $84.0\%$ charge
$\rightarrow$ grid-gated censor $84.7\%$ / $77.9\%$ $\rightarrow$ full censor $83.9\%$ / $75.6\%$.
\textbf{Why it fails, two mechanisms:} (1) classic deconvolution scores an envelope using the same
neutral mass at \emph{other} charge states (\texttt{observe\_adjacent\_charge\_states}); those
cross-charge peaks are off this feature's \mz\ grid, so censoring deletes the very cross-charge
evidence decon relies on --- which is why the \emph{charge}-recall column drops hardest ($-6\%$).
(2) The averaging normalises per-scan by TIC, so removing peaks reshuffles the composite. Combined
with the noisy claim partition (Section~\ref{sec:detector}), censoring deletes real teeth.
\textbf{Lesson:} already-assigned peaks are not merely interferents --- deconvolution exploits
neighbouring and cross-charge peaks, so a claim partition must not be read as a segmentation.

\subsection{Coverage depth as a recovery lever}
\label{sec:neg-coverage}
Raising the coverage target to detect more features ($0.90 \rightarrow 0.95$;
$\sim\!\num{96000} \rightarrow \num{143000}$ features) recovered \emph{zero} additional reference
peaks and left the miss set byte-identical (Section~\ref{sec:halving}). \textbf{Why it fails:} the
misses were not undetected --- they were detected and then mislabelled (charge, mono, or knitting).
Detection depth is the wrong knob when the loss is downstream; more darts do not help when the
darts already land and are then mis-scored.

\section{Synthesis: design principles}
\label{sec:synthesis}

The failures above, though they surface at different stages, teach a small set of transferable
principles for isotope-envelope feature detectors.

\begin{enumerate}[leftmargin=1.6em,label=\textbf{\arabic*.}]
  \item \textbf{A scoring metric must see signal on both sides of its hypothesis.} A one-sided
  window (mono $-0.5\Delta$ upward) is structurally blind to an error in the excluded direction
  (mono-too-high). Either make the metric two-sided or add an explicit, tightly-gated re-scoring
  against the alternatives whose windows reach the missing evidence (Sections~\ref{sec:walkback},
  \ref{sec:neg-window}).
  \item \textbf{Harmonic ambiguity needs absolute quality gates, not just relative comparison.}
  The $z \leftrightarrow 2z$ degeneracy is inherent to comb matching. Because the true hypothesis can
  itself fit poorly in crowded windows, a relative margin cannot protect it; combine loyalty to an
  independent prior with an \emph{absolute} fit floor so a charge is overridden only by an
  alternative that is good in absolute terms (Sections~\ref{sec:halving}, \ref{sec:doubling}).
  \item \textbf{Chromatographic co-elution is a discriminator the mass domain lacks.} Species
  $\sim 1\Da$ apart are indistinguishable to an off-by-one mass test but separable by elution
  profile. Gate mass-based knitting on mutual apex-in-window co-elution (Section~\ref{sec:knit}).
  \item \textbf{Exploit cross-charge redundancy, weighted by quality.} A peptide's other charge
  states are the richest resource for resolving an ambiguous, contaminated, or mono-missing charge
  --- but only if consensus weights fit quality, with independent multi-charge \emph{support} as the
  primary key so neither a loud contaminated charge nor a pair of concordant-but-wrong siblings can
  override a clean placement (Sections~\ref{sec:vote}, \ref{sec:neg-fitprimary}).
  \item \textbf{Anchor estimation on the most robust observable.} Anchor deconvolution on the
  highest-SNR peak (or an upstream committed claim), never on the fragile quantity being estimated;
  this prevents both the distant grab and circular mono-keying (Section~\ref{sec:anchor}).
  \item \textbf{Distinguish algorithm errors from data-limited inherent ambiguity.} Some misses are
  genuine data gaps (mono absent at a charge) or questionable references; the correct response is
  architectural (trust the charge that sees clean signal), not an ever-more-aggressive heuristic that
  regresses recoverable cases (Section~\ref{sec:datagaps}).
  \item \textbf{Adversarial and negative experiments are load-bearing.} Four of the most tempting
  ``obvious'' improvements (wider window, fit-primary consensus, broad walk-back, censoring) all
  regressed recall, each for an instructive reason. Measuring and reverting them --- and recording
  \emph{why} --- is what prevents the same intuitions from being re-implemented (Section~\ref{sec:negative}).
  \item \textbf{Fix an error where the richest evidence lives.} The charge-halving error originates
  in the detector but is best fixed in refine, where a high-SNR apex composite and a
  completeness-aware metric are both available (Sections~\ref{sec:halving}, \ref{sec:detector}).
\end{enumerate}

\section{Open problems and remaining misses}
\label{sec:open}

After the fixes above, $\sim 18$ reference peaks ($602/620$) remain unrecovered (an earlier snapshot:
$22$ misses = $3$ genuinely absent $+ 19$ off-by-one). Off-by-one remains the single largest miss
category. The open problems:

\begin{itemize}[leftmargin=1.4em]
  \item \textbf{Contaminated single-charge features with no clean sibling.} The cross-charge
  backstop (Section~\ref{sec:vote}) needs a second charge that sees the clean signal. When a peptide
  is observed at only one charge and that charge's envelope is contaminated (low absolute fit, e.g.
  $\sim0.25$--$0.39$), there is no clean vote to prefer and relative fit is unreliable. A
  chimera-robust single-envelope discriminator (or a genuine multi-envelope NNLS unmixing of the
  crowded window) is the plausible next step; a naive cosine corrector for this case regressed recall
  in earlier work.
  \item \textbf{Data-dependent averaging-window sizing.} The composite is hard-capped at $3$ scans
  (apex $\pm 1$), tuned to the CA/Lumos $10$-min gradient's narrow $\sim1.9$~s FWHM. This is not
  appropriate for longer gradients, where a feature spans more MS1 scans and averaging only $3$ would
  discard real signal and clip the elution. The constant should be replaced by a scan count derived
  from the measured chromatographic FWHM (the detector already measures FWHM via XIC half-max for its
  RT $\sigma$), so it adapts to gradient length instead of being hard-coded.
  \item \textbf{Genuine data gaps and reference quality.} Some remaining ``misses'' are cases where
  the monoisotope is absent from the centroids at every observed charge, or where the reference peak
  itself is questionable. These bound the achievable recall and should be characterised (via the
  miss-explorer envelopes) rather than optimised against, to avoid regressing recoverable cases.
  \item \textbf{Precision.} The shift-apex default roughly doubled the resolved-feature count versus
  the classic path (the classic \texttt{None}-drop was an implicit quality filter). Recall improved,
  but a precision/false-discovery accounting against the reference is an outstanding check.
\end{itemize}

\section*{Summary tables}
\addcontentsline{toc}{section}{Summary tables}

\begin{table}[h]
\centering
\small
\renewcommand{\arraystretch}{1.25}
\begin{tabular}{@{}p{2.0cm} p{3.5cm} p{3.2cm} p{4.4cm}@{}}
\toprule
\textbf{Stage} & \textbf{Failure / mechanism} & \textbf{Example peptide} & \textbf{Fix} \\
\midrule
Refine (mono) & Heavy off-by-one-too-high; fit-window blind below mono &
  ECCHGDLLECADDRADLAK (\zeq{4}, $2246.94\Da$) &
  Mono walk-back, gated mass $\geq 2200\Da$ \& $|z|\geq 3$ \\
Detect / Refine (charge) & Charge-halving: even \zeq{2} isotopes form a valid \zeq{1} comb &
  EFLDVIK (\zeq{2}, $862.48\Da$) &
  Envelope-fit recharge over $\{z/2,z,2z\}$ \\
Refine (charge) & Charge-doubling: \zeq{z} teeth are a subset of \zeq{2z}; interferents fill gaps &
  HVGDLGNVTAD[Na]K (\zeq{2}, $1246.59\Da$) &
  Detector-charge loyalty $+$ absolute fit floor $0.85$ \\
Refine (anchor) & Distant grab: anchor on stronger foreign peak in wide window &
  (four-way comparator, dataset-wide) &
  Anchor on detector's own claimed peak \\
Resolve (knit) & Separately-eluting species merged as false off-by-one &
  ECC \zeq{4} vs \zeq{2}/\zeq{3} interloper &
  Co-elution gate on off-by-one ($k\neq0$) links \\
Resolve (vote) & Consensus by intensity/count over-rides clean placement &
  AAVTAFWGK (\zeq{1} fit $0.39$ vs \zeq{2} $0.99$) &
  Fit as tiebreak \emph{after} support (not primary) \\
Data-limited & Mono absent from centroids at some charges &
  ECC at \zeq{2}/\zeq{3} &
  Trust the charge that sees the mono (\zeq{4}) \\
Detect & \texttt{RawSum} ignores missing teeth \& unexplained peaks &
  (root of charge-halving class) &
  Envelope-fit cosine downstream \\
\bottomrule
\end{tabular}
\caption{Failure-mode taxonomy: stage, root mechanism, traced example, and fix.}
\label{tab:summary}
\end{table}

\begin{table}[h]
\centering
\renewcommand{\arraystretch}{1.25}
\begin{tabular}{@{}l l c c@{}}
\toprule
\textbf{Step} & \textbf{Fix applied} & \textbf{Recall / 620} & \textbf{Charge-matched} \\
\midrule
Baseline (classic decon) & --- & $555$ ($89.5\%$) & $529$ ($85.3\%$) \\
Detector-anchored shift-apex $+$ recharge & Secs.~\ref{sec:halving},~\ref{sec:distant-grab} & $589$ ($95.0\%$) & $564$ ($91.9\%$) \\
Mono walk-back & Sec.~\ref{sec:walkback} & $589$ & $566$ \\
Co-elution knitting gate & Sec.~\ref{sec:knit} & $598$ ($96.5\%$) & $570$ \\
Fit tiebreak in consensus & Sec.~\ref{sec:vote} & $599$ ($96.6\%$) & $570$--$573$ \\
Anti-doubling fit floor & Sec.~\ref{sec:doubling} & $602$ ($97.1\%$) & $584$ ($94.2\%$) \\
\bottomrule
\end{tabular}
\caption{Recall progression across the session's fixes on CA/Lumos cov90 vs FlashLFQ. Intermediate
charge-matched values ($566$, $570$, $573$) are reported where the source commits recorded them; the
$589 \rightarrow 602$ recall trajectory and $529/564/584$ charge endpoints are the load-bearing
figures. The $89.5\% \rightarrow 95.0\%$ jump is the shift-apex$+$recharge default; the walk-back
alone left \emph{recall} flat at $589$ while lifting charge to $566$ (it corrects mass placement on
already-counted features), with the co-elution gate delivering the next $+9$.}
\label{tab:progression}
\end{table}

\section*{Sources and reproducibility}
\addcontentsline{toc}{section}{Sources and reproducibility}
All numbers are drawn from the pipeline's commit history and verified where possible against the
data files. The reference peptide masses, charges, RTs, and \mz{} values in
Sections~\ref{sec:mono}--\ref{sec:datagaps} were re-read directly from
\texttt{AllQuantifiedPeaks.tsv} (wide-tolerance search) and agree with the traced values: ECC \zeq{4}
$2246.9355\Da$ RT $12.153$ \mz\ $562.997$; HVG[Na] \zeq{2} $1246.5918\Da$ RT $11.238$ \mz\ $624.309$
intensity $2.75\times10^{6}$; AAVTAFWGK \zeq{1} $992.508\Da$ RT $16.31$ \mz\ $993.525$; DSC \zeq{3}
$3183.3808\Da$ RT $15.414$ \mz\ $1062.478$; EFLDVIK \zeq{2} theoretical $862.480\Da$ RT $14.70$ \mz\
$432.251$. The full envelope-fit metric derivation (including the $S^2 = R^2$ explained-energy
identity and the EFLDVIK completeness/explained-fraction table) is in the companion note
\emph{An Envelope-Fit Metric for Charge-State and Monoisotope Selection}
(\texttt{agent\_info/Envelope-Fit-Metric.tex}). Fit scores quoted for HVG (\zeq{1} $0.235$, \zeq{2}
$0.258$, \zeq{4} $0.689$), AAVTAFWGK (\zeq{1} $0.39$, \zeq{2} $0.99$), and the walk-back ($0.934$ at
$+1$) are from the traced-run diagnostics recorded in the commit messages and were not independently
re-derived here.

\end{document}
